\SourceLeaf{152}
Anmerkung in den A. N. Nr. 110, p. 230\HistoricalNote{[*)]}{[*)}{Vgl. Band VI, S. 457.]}, mit im Sinn gehabt —); aber darin hat er gefehlt, dass er den mittlern zu befürchtenden Fehler in den Resultaten nicht mit berechnet hat, und dass er durch einen Umstand, welchen Herr \textsc{Ivory} mit Recht tadelt, nemlich durch die Zusammenstellung der verschiedenen Resultate aus verschiedenen Combinationen, wobei immer seine Beobachtungen in der Nähe des Äquators einen überwiegenden Einfluss behalten müssen, den Schein einer viel grössern Genauigkeit der Endresultate hervorbringt, als diese wirklich haben. In der That ist es damit gewissermassen so, obwohl lange nicht in dem Masse, wie der verstorbene \textsc{Ende} zuweilen Polhöhenbestimmungen nach \textsc{Douwes}’ Methode machte, indem er immer \textls{dieselbe} Höhe in der Nähe des Mittags mit \textls{verschiedenen} weit davon entfernten verband und dann Resultate fand, die innerhalb $0'',01$ übereinstimmten.

Einzeln widerlegen kann man diesen Aufsatz von \textsc{Ivory} eben so schwer, da darin eben so wenig logische Ordnung ist; ich finde auch nicht den allerkleinsten Grund gegen die Zweckmässigkeit der M. der kl. Q. darin, aber was soll man von dem Verfahren sagen, welches \textsc{Ivory} dafür substituiren will, indem er die kleinste Pendellänge mit vielen in grosser Breite combinirt, das ist ja eines verständigen Mathematikers ganz unwürdig und eigentlich das \textsc{Ende}sche Verfahren\HistoricalNote{[**)]}{[**)}{Die vorstehende Kritik bezieht sich auf folgende Arbeiten J. \textsc{Ivorys}:
\par\hangindent1.5em 1. On the method of least squares. Phil. Mag. (Tilloch), \textsc{lxviii}, 1826, S. 161—165.
\par\hangindent1.5em 2. On the Ellipticity of the Earth, as deduced from Experiments made with the Pendulum. Phil. Mag. (Tilloch), \textsc{lxviii}, 1826, S. 3—10, 92—101, 241—245, 246—251, 321—326, 350—353.]}. . . . . . . .

\Needspace{6\baselineskip}
\begin{center}\textsc{Gauss} an \textsc{Encke}. Göttingen, 23. August 1831.\end{center}

. . . . . . . . . . . . . .

Nicht ohne Interesse habe ich aus Ihrem Briefe den Gang gesehen, den Sie zur Rechtfertigung des Verfahrens, das arithmetische Mittel zu nehmen, eingeschlagen haben. Ich finde diesen Gang sehr beifallswerth, in so fern auf die Frage, was zu thun sei, eine von allen Betrachtungen der Wahrscheinlich%
